%% ============================================================
%%  SECCIÓN 5 — Momento Angular Tensorial, Álgebra so(3) y CCOC
%% ============================================================

\section{Momento Angular Tensorial, Álgebra de Lie \texorpdfstring{$\mathfrak{so}(3)$}{so(3)} y el CCOC}
\label{sec:momento_angular}

\subsection{Herramientas Tensoriales en \texorpdfstring{$\mathbb{R}^3$}{R3}: Notación de Índices, Kronecker y Levi-Civita}

Para abordar el cálculo de conmutadores en tres dimensiones espaciales sin recurrir a cálculos algebraicos interminables, la física teórica moderna emplea la \textbf{notación de índices} y el convenio de suma de Einstein \cite{sakurai, zettili}:
\begin{itemize}
    \item Los índices latinos $i, j, k, l, m, \dots \in \{1, 2, 3\}$ denotan las componentes espaciales cartesianas $(x, y, z) \equiv (x_1, x_2, x_3)$.
    \item \textbf{Convenio de Suma de Einstein:} Todo índice repetido en un mismo término monomial implica una sumatoria sobre sus tres valores posibles:
    \begin{equation}
        A_i B_i \equiv \sum_{i=1}^3 A_i B_i = A_1 B_1 + A_2 B_2 + A_3 B_3 = \vec{A}\cdot\vec{B}.
    \end{equation}
\end{itemize}

\subsubsection{El Delta de Kronecker \texorpdfstring{$\delta_{ij}$}{delta\_ij}}
Se define como el tensor de segundo orden:
\begin{equation}
    \delta_{ij} \equiv \begin{cases} 1, & \text{si } i = j, \\ 0, & \text{si } i \neq j. \end{cases}
\end{equation}
Actúa como operador de sustitución de índices: $\delta_{ij} A_j = A_i$, y su traza en 3D es $\delta_{ii} = \delta_{11} + \delta_{22} + \delta_{33} = 3$.

\subsubsection{El Pseudotensor Totalmente Antisimétrico de Levi-Civita \texorpdfstring{$\varepsilon_{ijk}$}{epsilon\_ijk}}
Se define como el tensor de permutación en tres dimensiones:
\begin{equation}
    \varepsilon_{ijk} \equiv \begin{cases}
        +1, & \text{si } (i, j, k) \text{ es una permutación par de } (1, 2, 3), \\
        -1, & \text{si } (i, j, k) \text{ es una permutación impar de } (1, 2, 3), \\
        0,  & \text{si cualesquiera dos índices son iguales}.
    \end{cases}
\end{equation}
Permite expresar el producto vectorial euclídeo de dos vectores $\vec{A} \times \vec{B}$ mediante sus componentes:
\begin{equation}
    (\vec{A} \times \vec{B})_i = \varepsilon_{ijk} A_j B_k.
\end{equation}

\begin{teoremabox}[Identidad de Contracción Fundamental de Levi-Civita]
    El producto de dos tensores de Levi-Civita con un índice contraído (sumado) se expresa en términos de deltas de Kronecker como:
    \begin{equation}
        \boxed{\varepsilon_{ijk}\varepsilon_{lmk} = \delta_{il}\delta_{jm} - \delta_{im}\delta_{jl}}.
        \label{eq:levi_civita_contract_1}
    \end{equation}
    Si además contraemos un segundo índice ($m = j$):
    \begin{equation}
        \varepsilon_{ijk}\varepsilon_{ljk} = \delta_{il}\delta_{jj} - \delta_{ij}\delta_{jl} = 3\delta_{il} - \delta_{il} = \boxed{2\delta_{il}}.
        \label{eq:levi_civita_contract_2}
    \end{equation}
\end{teoremabox}

\begin{alertabox}[¿De dónde surge el factor 2 en la contracción de dos índices?]
    Para comprender físicamente la Ec.~\eqref{eq:levi_civita_contract_2}, consideremos fijar $i = 1$ y $l = 1$. La suma sobre $j$ y $k$ recorre:
    \begin{equation}
        \sum_{j,k=1}^3 \varepsilon_{1jk}\varepsilon_{1jk} = \varepsilon_{123}\varepsilon_{123} + \varepsilon_{132}\varepsilon_{132} = (+1)(+1) + (-1)(-1) = 1 + 1 = 2.
    \end{equation}
    Los únicos pares de índices que no se anulan son aquellos donde $j$ y $k$ son diferentes entre sí y distintos de $i=1$. Existen exactamente $2! = 2$ permutaciones posibles de dichos índices restantes, cada una contribuyendo con $(+1)^2 = 1$ o $(-1)^2 = 1$. Por ello, el factor combinatorio resultante es exactamente $2$.
\end{alertabox}

\subsection{Definición Cuántica del Operador Momento Angular}

En mecánica clásica, el momento angular de una partícula puntual respecto al origen es $\vec{L} = \vec{r} \times \vec{p}$. Por el principio de correspondencia de Dirac, reemplazamos las coordenadas y momentos canónicos por operadores hermíticos:
\begin{equation}
    [\hat{x}_i, \hat{p}_j] = \iu\hbar\delta_{ij}\hat{I}, \quad [\hat{x}_i, \hat{x}_j] = 0, \quad [\hat{p}_i, \hat{p}_j] = 0.
    \label{eq:canonical_comm_3d}
\end{equation}
El operador momento angular orbital cuántico se define en componentes tensoriales por:
\begin{equation}
    \boxed{\hat{L}_k \equiv \varepsilon_{kij}\hat{x}_i\hat{p}_j}.
    \label{eq:l_operator_def}
\end{equation}
Desarrollando explícitamente para cada componente cartesiana:
\begin{align}
    \hat{L}_1 &= \hat{L}_x = \varepsilon_{123}\hat{x}_2\hat{p}_3 + \varepsilon_{132}\hat{x}_3\hat{p}_2 = \hat{x}_2\hat{p}_3 - \hat{x}_3\hat{p}_2 = \hat{y}\hat{p}_z - \hat{z}\hat{p}_y, \\
    \hat{L}_2 &= \hat{L}_y = \varepsilon_{231}\hat{x}_3\hat{p}_1 + \varepsilon_{213}\hat{x}_1\hat{p}_3 = \hat{x}_3\hat{p}_1 - \hat{x}_1\hat{p}_3 = \hat{z}\hat{p}_x - \hat{x}\hat{p}_z, \\
    \hat{L}_3 &= \hat{L}_z = \varepsilon_{312}\hat{x}_1\hat{p}_2 + \varepsilon_{321}\hat{x}_2\hat{p}_1 = \hat{x}_1\hat{p}_2 - \hat{x}_2\hat{p}_1 = \hat{x}\hat{p}_y - \hat{y}\hat{p}_x.
\end{align}
Notemos que como $i \neq j$ en cada producto $\hat{x}_i\hat{p}_j$, las componentes de posición y momento involucradas conmutan entre sí ($[\hat{x}_i, \hat{p}_j] = \iu\hbar\delta_{ij} = 0$ para $i \neq j$), lo que garantiza que $\hat{L}_k$ es un operador estrictamente \textbf{hermítico}:
\begin{equation}
    \hat{L}_k\adj = (\varepsilon_{kij}\hat{x}_i\hat{p}_j)\adj = \varepsilon_{kij}\hat{p}_j\adj\hat{x}_i\adj = \varepsilon_{kij}\hat{p}_j\hat{x}_i = \varepsilon_{kij}\hat{x}_i\hat{p}_j = \hat{L}_k.
\end{equation}

\subsection{Demostración Rigurosa del Álgebra de Lie \texorpdfstring{$\mathfrak{so}(3)$}{so(3)}}

Deseamos calcular el conmutador fundamental entre dos componentes arbitrarias $[\hat{L}_i, \hat{L}_j]$ sin recurrir al cálculo manual de las tres componentes por separado.

\begin{teoremabox}[Álgebra de Lie del Momento Angular]
    Las componentes del operador momento angular cuántico satisfacen las relaciones de conmutación:
    \begin{equation}
        \boxed{[\hat{L}_i, \hat{L}_j] = \iu\hbar\,\varepsilon_{ijk}\hat{L}_k}.
        \label{eq:so3_algebra}
    \end{equation}
\end{teoremabox}

\subsubsection{Del Cálculo Cartesiano Manual a la Deducción Tensorial General}

\begin{alertabox}[{Calentamiento: Cálculo Cartesiano Manual de $[\hat{L}_x, \hat{L}_y]$ frente al Método Tensorial}]
    Antes de efectuar la demostración general con índices para las nueve combinaciones $[\hat{L}_i, \hat{L}_j]$, evaluemos manualmente un caso particular: $[\hat{L}_x, \hat{L}_y] = [\hat{y}\hat{p}_z - \hat{z}\hat{p}_y, \, \hat{z}\hat{p}_x - \hat{x}\hat{p}_z]$.\\
    Por bilinealidad, este conmutador se expande en cuatro términos:
    \begin{equation}
        [\hat{L}_x, \hat{L}_y] = [\hat{y}\hat{p}_z, \hat{z}\hat{p}_x] - \bcancel{[\hat{y}\hat{p}_z, \hat{x}\hat{p}_z]} - \bcancel{[\hat{z}\hat{p}_y, \hat{z}\hat{p}_x]} + [\hat{z}\hat{p}_y, \hat{x}\hat{p}_z].
    \end{equation}
    Los dos términos centrales se anulan porque todas las coordenadas y momentos que contienen pertenecen a ejes distintos y conmutan entre sí. En el primer y cuarto término, la única pareja que no conmuta es $\hat{z}$ con $\hat{p}_z$ ($[\hat{p}_z, \hat{z}] = -\iu\hbar\hat{I}$ y $[\hat{z}, \hat{p}_z] = +\iu\hbar\hat{I}$), mientras que $\hat{x}, \hat{p}_x, \hat{y}, \hat{p}_y$ salen como factores espectadores:
    \begin{equation}
        [\hat{L}_x, \hat{L}_y] = \hat{y}\hat{p}_x [\hat{p}_z, \hat{z}] + \hat{x}\hat{p}_y [\hat{z}, \hat{p}_z] = \hat{y}\hat{p}_x(-\iu\hbar\hat{I}) + \hat{x}\hat{p}_y(+\iu\hbar\hat{I}) = \iu\hbar(\hat{x}\hat{p}_y - \hat{y}\hat{p}_x) = \iu\hbar\hat{L}_z.
    \end{equation}
    \textbf{¿Qué hace el método tensorial que veremos a continuación?} Exactamente esta misma operación para los 9 pares $(i, j)$ simultáneamente. Para no confundir la doble sumatoria interna de $\hat{L}_i$ con la doble sumatoria interna de $\hat{L}_j$, bautizamos con \textbf{nombres de índices mudos distintos} a cada una: $(a, b)$ para $\hat{L}_i = \varepsilon_{iab}\hat{x}_a\hat{p}_b$ y $(c, d)$ para $\hat{L}_j = \varepsilon_{jcd}\hat{x}_c\hat{p}_d$.
\end{alertabox}

\textbf{Demostración Tensorial Paso a Paso:}\\
Insertamos la definición tensorial \eqref{eq:l_operator_def} empleando índices mudos independientes para no mezclar las sumas:
\begin{equation}
    [\hat{L}_i, \hat{L}_j] = \left[\varepsilon_{iab}\hat{x}_a\hat{p}_b, \, \varepsilon_{jcd}\hat{x}_c\hat{p}_d\right] = \varepsilon_{iab}\varepsilon_{jcd}\left[\hat{x}_a\hat{p}_b, \, \hat{x}_c\hat{p}_d\right].
    \label{eq:comm_ll_tensor}
\end{equation}
Aplicamos la identidad de Leibniz para el conmutador de dos productos de operadores:
\begin{equation}
    [\hat{A}\hat{B}, \hat{C}\hat{D}] = \hat{A}[\hat{B}, \hat{C}]\hat{D} + \hat{A}\hat{C}[\hat{B}, \hat{D}] + [\hat{A}, \hat{C}]\hat{D}\hat{B} + \hat{C}[\hat{A}, \hat{D}]\hat{B}.
\end{equation}
Evaluando para nuestras variables canónicas, y recordando que $[\hat{x}_a, \hat{x}_c] = 0$ y $[\hat{p}_b, \hat{p}_d] = 0$:
\begin{align}
    [\hat{x}_a\hat{p}_b, \, \hat{x}_c\hat{p}_d] &= \hat{x}_a [\hat{p}_b, \hat{x}_c]\hat{p}_d + \hat{x}_c [\hat{x}_a, \hat{p}_d]\hat{p}_b \nonumber\\
    &= \hat{x}_a(-\iu\hbar\delta_{bc})\hat{p}_d + \hat{x}_c(\iu\hbar\delta_{ad})\hat{p}_b \nonumber\\
    &= -\iu\hbar\delta_{bc}\hat{x}_a\hat{p}_d + \iu\hbar\delta_{ad}\hat{x}_c\hat{p}_b.
    \label{eq:leibniz_expansion}
\end{align}
Sustituyendo esta expresión dentro de la Ec.~\eqref{eq:comm_ll_tensor}, el conmutador se separa en dos términos: $[\hat{L}_i, \hat{L}_j] = \mathcal{T}_1 + \mathcal{T}_2$.

\textit{Evaluación del Primer Término $\mathcal{T}_1$:}
\begin{align}
    \mathcal{T}_1 &= \varepsilon_{iab}\varepsilon_{jcd}(-\iu\hbar\delta_{bc})\hat{x}_a\hat{p}_d = -\iu\hbar\,\varepsilon_{iab}\varepsilon_{jbd}\hat{x}_a\hat{p}_d.
\end{align}
Invertimos los dos últimos índices del segundo Levi-Civita para alinearlo con la identidad de contracción ($\varepsilon_{jbd} = -\varepsilon_{jdb}$):
\begin{equation}
    \mathcal{T}_1 = \iu\hbar\,\varepsilon_{iab}\varepsilon_{jdb}\hat{x}_a\hat{p}_d.
\end{equation}
Contraemos sobre el índice repetido $b$ aplicando la identidad \eqref{eq:levi_civita_contract_1}:
\begin{equation}
    \varepsilon_{iab}\varepsilon_{jdb} = \delta_{ij}\delta_{ad} - \delta_{id}\delta_{aj}.
\end{equation}
Por ende:
\begin{align}
    \mathcal{T}_1 &= \iu\hbar\left(\delta_{ij}\delta_{ad} - \delta_{id}\delta_{aj}\right)\hat{x}_a\hat{p}_d \nonumber\\
    &= \iu\hbar\,\delta_{ij}(\hat{x}_a\hat{p}_a) - \iu\hbar\,\hat{x}_j\hat{p}_i.
    \label{eq:t1_eval}
\end{align}

\textit{Evaluación del Segundo Término $\mathcal{T}_2$:}
\begin{align}
    \mathcal{T}_2 &= \varepsilon_{iab}\varepsilon_{jcd}(\iu\hbar\delta_{ad})\hat{x}_c\hat{p}_b = \iu\hbar\,\varepsilon_{iab}\varepsilon_{jca}\hat{x}_c\hat{p}_b.
\end{align}
Reordenamos cíclicamente los índices para contraer sobre $a$: $\varepsilon_{iab} = \varepsilon_{abi}$, y $\varepsilon_{jca} = \varepsilon_{ajc} = -\varepsilon_{acj}$:
\begin{equation}
    \mathcal{T}_2 = -\iu\hbar\,\varepsilon_{abi}\varepsilon_{acj}\hat{x}_c\hat{p}_b.
\end{equation}
Aplicando la identidad de contracción sobre $a$:
\begin{equation}
    \varepsilon_{abi}\varepsilon_{acj} = \delta_{bc}\delta_{ij} - \delta_{bj}\delta_{ic}.
\end{equation}
Por ende:
\begin{align}
    \mathcal{T}_2 &= -\iu\hbar\left(\delta_{bc}\delta_{ij} - \delta_{bj}\delta_{ic}\right)\hat{x}_c\hat{p}_b \nonumber\\
    &= -\iu\hbar\,\delta_{ij}(\hat{x}_b\hat{p}_b) + \iu\hbar\,\hat{x}_i\hat{p}_j.
    \label{eq:t2_eval}
\end{align}

\textit{Suma de Términos $\mathcal{T}_1 + \mathcal{T}_2$:}
Sumamos miembro a miembro las Ecs.~\eqref{eq:t1_eval} y \eqref{eq:t2_eval}:
\begin{equation}
    [\hat{L}_i, \hat{L}_j] = \iu\hbar\,\delta_{ij}(\hat{x}_a\hat{p}_a) - \iu\hbar\,\hat{x}_j\hat{p}_i - \iu\hbar\,\delta_{ij}(\hat{x}_b\hat{p}_b) + \iu\hbar\,\hat{x}_i\hat{p}_j.
\end{equation}
Como los índices mudos $a$ y $b$ representan la misma suma $\vec{x}\cdot\vec{p}$, los términos diagonales proporcionales a $\delta_{ij}$ \textbf{se cancelan de forma idéntica}:
\begin{equation}
    [\hat{L}_i, \hat{L}_j] = \iu\hbar\left(\hat{x}_i\hat{p}_j - \hat{x}_j\hat{p}_i\right).
    \label{eq:l_comm_step}
\end{equation}
Reconocemos la definición tensorial del momento angular: $\varepsilon_{ijk}\hat{L}_k = \varepsilon_{ijk}\varepsilon_{kab}\hat{x}_a\hat{p}_b = (\delta_{ia}\delta_{jb} - \delta_{ib}\delta_{ja})\hat{x}_a\hat{p}_b = \hat{x}_i\hat{p}_j - \hat{x}_j\hat{p}_i$.
Sustituyendo esto en la Ec.~\eqref{eq:l_comm_step}:
\begin{equation}
    [\hat{L}_i, \hat{L}_j] = \iu\hbar\,\varepsilon_{ijk}\hat{L}_k.
\end{equation}
La deducción es exacta, rigurosa y general para cualquier par de componentes.

\subsubsection{Conexión con Álgebra Lineal: El Grupo $\mathrm{SO}(3)$ y Representación Hermítica}

\begin{alertabox}[Conexión con Álgebra Lineal: El Grupo $\mathrm{SO}(3)$ y los Generadores de Rotación]
    En Álgebra Lineal II y Geometría Afín, toda rotación en $\mathbb{R}^3$ alrededor de un eje viene dada por una matriz ortogonal $R \in \mathrm{SO}(3)$ ($\det R = +1$, $R^T R = I$), la cual puede escribirse como una exponencial matricial $R = \eu^{\theta \Omega}$, donde $\Omega$ es una matriz real \textbf{antisimétrica} ($\Omega^T = -\Omega$).
    
    El espacio vectorial de matrices antisimétricas de $3\times 3$ conforma el álgebra de Lie $\mathfrak{so}(3)$, cuya base canónica está formada por las tres matrices generadoras $J_1, J_2, J_3$ definidas por componentes mediante el pseudotensor de Levi-Civita:
    \begin{equation}
        (J_i)_{jk} \equiv -\varepsilon_{ijk}.
    \end{equation}
    Un cálculo directo de conmutadores matriciales arroja el álgebra geométrica clásica:
    \begin{equation}
        [J_i, J_j] = \varepsilon_{ijk} J_k.
    \end{equation}
    En mecánica cuántica, los observables físicos deben ser operadores autoadjuntos (hermíticos, $\hat{L}_i\adj = \hat{L}_i$). Como las matrices geométricas $J_i$ son antisimétricas reales ($J_i\adj = J_i^T = -J_i$), multiplicarlas por $+\iu\hbar$ las convierte rigurosamente en operadores hermíticos:
    \begin{equation}
        \hat{L}_i \equiv +\iu\hbar\,J_i \implies \hat{L}_i\adj = (-\iu\hbar)(J_i\adj) = (-\iu\hbar)(-J_i) = +\iu\hbar\,J_i = \hat{L}_i.
    \end{equation}
    \textbf{¿Por qué multiplicamos exactamente por $+\iu\hbar$? (Intuición Física y Álgebra Lineal):}
    \begin{itemize}
        \item \textbf{El factor imaginario $\iu$:} Toda matriz antisimétrica real posee autovalores puramente imaginarios ($\lambda \in \iu\mathbb{R}$). Al multiplicar por $\iu$, dichos autovalores se convierten en números estrictamente \textbf{reales} ($\iu \cdot \iu\lambda = -\lambda \in \mathbb{R}$), satisfaciendo la exigencia cuántica de que toda medición experimental en el laboratorio arroje números reales.
        \item \textbf{La constante de Planck $\hbar$:} Las matrices geométricas $J_i$ son magnitudes adimensionales (ángulos puros en radianes), mientras que el momento angular físico posee dimensiones de acción ($\mathrm{J}\cdot\mathrm{s}$). La constante $\hbar$ dota al generador de las unidades dinámicas del universo cuántico.
    \end{itemize}
    Multiplicando el conmutador clásico $[J_i, J_j] = \varepsilon_{ijk}J_k$ por $(\iu\hbar)^2 = -\hbar^2$:
    \begin{align}
        [\hat{L}_i, \hat{L}_j] &= [\iu\hbar\,J_i, \, \iu\hbar\,J_j] = (\iu\hbar)^2 [J_i, J_j] = (-\hbar^2)\varepsilon_{ijk}J_k \nonumber\\
        &= \iu\hbar\,\varepsilon_{ijk}(\iu\hbar\,J_k) = \iu\hbar\,\varepsilon_{ijk}\hat{L}_k.
    \end{align}
    Esta identificación es además directamente consistente con el operador de rotación finita: si escribimos la rotación espacial como $R = \eu^{\theta J}$ con $J = \frac{1}{\iu\hbar}\hat{L} = -\frac{\iu}{\hbar}\hat{L}$, recuperamos el operador unitario cuántico de rotación $R(\theta) = \eu^{-\iu\theta\hat{L}/\hbar}$.\\
    \textbf{Representación Vectorial:} Cabe subrayar que estas matrices $3\times 3$ constituyen la \textbf{representación vectorial o adjunta} (correspondiente al subespacio irreductible de momento angular $\ell=1$) del álgebra de Lie $\mathfrak{so}(3)$, a diferencia de los operadores diferenciales actuando sobre el espacio funcional de dimensión infinita $\mathcal{L}^2(\mathbb{R}^3)$, que representan la misma álgebra abstracta.\\
    ¡El álgebra cuántica del momento angular no es un postulado misterioso! Es la \textbf{representación hermítica del álgebra de rotaciones de $\mathbb{R}^3$} que se estudia en álgebra lineal.
\end{alertabox}

\subsection{La Paradoja del Producto Cruz Cuántico: \texorpdfstring{$\hat{\vec{L}} \times \hat{\vec{L}} = \iu\hbar\hat{\vec{L}} \neq \vec{0}$}{L x L = ihbar L}}

En el álgebra vectorial elemental de $\mathbb{R}^3$, el producto vectorial de cualquier vector consigo mismo es idénticamente nulo: $\vec{v} \times \vec{v} = \vec{0}$, debido a que $\sin(0^\circ) = 0$ y a la antisimetría del determinante.

¿Qué ocurre en el universo cuántico cuando evaluamos el producto vectorial del vector de operadores $\hat{\vec{L}}$ consigo mismo?
\begin{equation}
    (\hat{\vec{L}} \times \hat{\vec{L}})_k \equiv \varepsilon_{kij}\hat{L}_i\hat{L}_j.
\end{equation}
Dado que el índice $k$ está fijo, los índices $i$ y $j$ se suman sobre las parejas posibles. Utilizando el truco algebraico de intercambiar índices mudos:
\begin{equation}
    \varepsilon_{kij}\hat{L}_i\hat{L}_j = \frac{1}{2}\varepsilon_{kij}\hat{L}_i\hat{L}_j + \frac{1}{2}\varepsilon_{kji}\hat{L}_j\hat{L}_i.
\end{equation}
Por la antisimetría de Levi-Civita, $\varepsilon_{kji} = -\varepsilon_{kij}$:
\begin{equation}
    (\hat{\vec{L}} \times \hat{\vec{L}})_k = \frac{1}{2}\varepsilon_{kij}(\hat{L}_i\hat{L}_j - \hat{L}_j\hat{L}_i) = \frac{1}{2}\varepsilon_{kij}[\hat{L}_i, \hat{L}_j].
\end{equation}
¡El producto vectorial se expresa directamente en términos del conmutador de operadores! Sustituyendo el álgebra de Lie recién demostrada $[\hat{L}_i, \hat{L}_j] = \iu\hbar\,\varepsilon_{ijm}\hat{L}_m$:
\begin{align}
    (\hat{\vec{L}} \times \hat{\vec{L}})_k &= \frac{1}{2}\varepsilon_{kij}\left( \iu\hbar\,\varepsilon_{ijm}\hat{L}_m \right) = \frac{\iu\hbar}{2}\left( \varepsilon_{kij}\varepsilon_{mij} \right)\hat{L}_m.
\end{align}
Aplicando la contracción de dos índices \eqref{eq:levi_civita_contract_2}, $\varepsilon_{kij}\varepsilon_{mij} = 2\delta_{km}$:
\begin{equation}
    (\hat{\vec{L}} \times \hat{\vec{L}})_k = \frac{\iu\hbar}{2}\left( 2\delta_{km} \right)\hat{L}_m = \iu\hbar\,\delta_{km}\hat{L}_m = \iu\hbar\,\hat{L}_k.
\end{equation}
Reuniendo las tres componentes en notación vectorial tridimensional:
\begin{equation}
    \boxed{\hat{\vec{L}} \times \hat{\vec{L}} = \iu\hbar\,\hat{\vec{L}} \neq \vec{0}}.
    \label{eq:quantum_cross_product}
\end{equation}
Este resultado es una de las firmas más bellas y distintivas del carácter no conmutativo de la mecánica cuántica: el vector momento angular \enquote{no puede alinearse consigo mismo} en el sentido clásico, pues sus componentes cartesianas son incompatibles entre sí \cite{sakurai}.

\subsection{El Operador Casimir \texorpdfstring{$\hat{L}^2$}{L\textasciicircum 2}\ e Invariancia Rotacional}

Definimos el operador escalar del momento angular al cuadrado (denominado formalmente el \textbf{operador Casimir cuadrático} del grupo de rotaciones $\mathrm{SO}(3)$):
\begin{equation}
    \hat{L}^2 \equiv \hat{L}_x^2 + \hat{L}_y^2 + \hat{L}_z^2 = \hat{L}_j\hat{L}_j.
    \label{eq:casimir_def}
\end{equation}

\begin{teoremabox}[Conmutación del Operador Casimir]
    El operador $\hat{L}^2$ conmuta simultáneamente con cada una de las componentes del momento angular:
    \begin{equation}
        \boxed{[\hat{L}^2, \hat{L}_i] = 0, \quad \forall\,i \in \{1, 2, 3\}}.
        \label{eq:l2_comm}
    \end{equation}
\end{teoremabox}

\subsubsection{Conmutación del Operador Casimir y Derivación de Leibniz}

\textbf{Demostración:}\\
Calculemos el conmutador aplicando la regla de Leibniz para productos:
\begin{align}
    [\hat{L}^2, \hat{L}_i] &= [\hat{L}_j\hat{L}_j, \hat{L}_i] = \hat{L}_j[\hat{L}_j, \hat{L}_i] + [\hat{L}_j, \hat{L}_i]\hat{L}_j.
\end{align}
Sustituyendo $[\hat{L}_j, \hat{L}_i] = \iu\hbar\,\varepsilon_{jik}\hat{L}_k = -\iu\hbar\,\varepsilon_{ijk}\hat{L}_k$:
\begin{align}
    [\hat{L}^2, \hat{L}_i] &= \hat{L}_j(-\iu\hbar\,\varepsilon_{ijk}\hat{L}_k) + (-\iu\hbar\,\varepsilon_{ijk}\hat{L}_k)\hat{L}_j \nonumber\\
    &= -\iu\hbar\,\varepsilon_{ijk}\left( \hat{L}_j\hat{L}_k + \hat{L}_k\hat{L}_j \right) = -\iu\hbar\,\varepsilon_{ijk}\{\hat{L}_j, \hat{L}_k\}.
    \label{eq:casimir_step}
\end{align}

\subsubsection{Teorema de Anulación: Contracción Simétrico-Antisimétrico}

\begin{alertabox}[¿Por qué la contracción de un tensor simétrico con uno antisimétrico se anula idénticamente?]
    Examinemos con lupa el término resultante $\varepsilon_{ijk}\{\hat{L}_j, \hat{L}_k\}$ sumado sobre los índices mudos $j$ y $k$:\\
    1. El anticonmutador es una cantidad \textbf{estrictamente simétrica} ante el intercambio de índices:
    \begin{equation}
        S_{jk} \equiv \{\hat{L}_j, \hat{L}_k\} = \hat{L}_j\hat{L}_k + \hat{L}_k\hat{L}_j = \{\hat{L}_k, \hat{L}_j\} = S_{kj}.
    \end{equation}
    2. El tensor de Levi-Civita $\varepsilon_{ijk}$ es \textbf{estrictamente antisimétrico} ante el intercambio de $j$ y $k$:
    \begin{equation}
        \varepsilon_{ikj} = -\varepsilon_{ijk}.
    \end{equation}
    Sea la suma escalar $T_i \equiv \varepsilon_{ijk} S_{jk}$. Como $j$ y $k$ son índices mudos de sumatoria, podemos cambiarles el nombre sin alterar el valor de la suma ($j \to k$ y $k \to j$):
    \begin{equation}
        T_i = \varepsilon_{ikj} S_{kj}.
    \end{equation}
    Utilizando la simetría de $S$ y la antisimetría de $\varepsilon$:
    \begin{equation}
        T_i = (-\varepsilon_{ijk}) S_{jk} = -(\varepsilon_{ijk} S_{jk}) = -T_i.
    \end{equation}
    La única cantidad matemática que es igual a su propio negativo es el cero:
    \begin{equation}
        T_i = -T_i \implies 2T_i = 0 \implies \boxed{T_i = \varepsilon_{ijk}\{\hat{L}_j, \hat{L}_k\} \equiv 0}.
    \end{equation}
    Sustituyendo esto en la Ec.~\eqref{eq:casimir_step}, se concluye irrefutablemente que $[\hat{L}^2, \hat{L}_i] = 0$.
\end{alertabox}

\subsection{Conjunto Completo de Observables Compatibles (CCOC)}

\subsubsection{Incompatibilidad Cuántica y Relaciones de Incertidumbre Transversales}

En mecánica clásica, es posible medir simultáneamente las tres componentes del vector momento angular $(L_x, L_y, L_z)$ con precisión infinita. En mecánica cuántica, sin embargo:
\begin{equation}
    [\hat{L}_x, \hat{L}_y] = \iu\hbar\hat{L}_z \neq 0.
\end{equation}
Por el teorema general de incompatibilidad cuántica, para cualquier estado físico con momento angular no nulo ($\ell \ge 1$), \textbf{es imposible medir simultáneamente más de una componente del momento angular con precisión absoluta}. Si preparamos un autoestado de $\hat{L}_z$ con autovalor $m\hbar$, las componentes transversales $\hat{L}_x$ y $\hat{L}_y$ poseen incertidumbres intrínsecas gobernadas por la relación de Heisenberg:
\begin{equation}
    \Delta L_x \Delta L_y \ge \frac{\hbar}{2}|\ev{\hat{L}_z}| = \frac{\hbar^2 |m|}{2}.
\end{equation}
Nótese que para autoestados con $m = 0$ pertenecientes a multipletes con $\ell \ge 1$, si bien la cota inferior directa en la desigualdad se anula ($\ev{\hat{L}_z}=0$), las incertidumbres cuánticas transversales reales siguen siendo estrictamente positivas: $\Delta L_x = \Delta L_y = \hbar\sqrt{\ell(\ell+1)/2} > 0$ (como se demuestra rigurosamente en el Problema Resuelto~3). Solo para el estado $\ell = 0$ (donde $\hat{\vec{L}}\ket{0, 0} = \vec{0}$) las tres componentes son simultáneamente nítidas con $\Delta L_x = \Delta L_y = \Delta L_z = 0$.

\begin{figure}[H]
    \centering
    \begin{tikzpicture}[>=Stealth, font=\small, scale=1.0]
  \def\L{2.4495}  % sqrt(6)
  % esfera de radio |L|
  \draw[grisUNI!60,dashed] (0,0) circle (\L);
  % eje z
  \draw[->,grisUNI] (0,-3.1)--(0,3.1) node[above]{$L_z$};
  \foreach \m in {-2,-1,0,1,2}{
    \pgfmathsetmacro{\r}{sqrt(6-\m*\m)}
    % cono: circunferencia de precesión (vista en perspectiva)
    \draw[azulUNI!70,dashed] (0,\m) ellipse ({\r} and {0.22*\r});
    % vector L con L_z = m hbar
    \draw[->,very thick,azulUNI] (0,0)--(\r,\m);
    \fill[azulUNI] (\r,\m) circle (1.6pt);
    % marca en el eje
    \draw (-0.07,\m)--(0.07,\m);
    \node[anchor=east,font=\footnotesize] at (-0.1,\m) {$%
      \ifnum\m=0 0\else\ifnum\m=1 \hbar\else\ifnum\m=-1 -\hbar\else \m\hbar\fi\fi\fi$};
  }
  % anotaciones
  \node[anchor=west,font=\footnotesize,align=left] at (3.0,2.05)
       {$|\vec L|=\hbar\sqrt{\ell(\ell+1)}=\hbar\sqrt6\approx2{,}45\,\hbar$};
  \draw[->,grisUNI] (3.0,1.85)--(1.95,1.55);
  \node[anchor=west,font=\footnotesize,align=left] at (3.0,-2.0)
       {$L_x,\,L_y$ indeterminados:\\ cada vector gira sobre su cono};
  \draw[<->,rojoUNI] (-2.75,-2.0)--(-2.75,2.0);
  \node[rojoUNI,font=\footnotesize,rotate=90] at (-3.05,0)
       {$|L_z|\le 2\hbar < |\vec L|$};
\end{tikzpicture}

    \caption{Modelo vectorial para $\ell=2$: los cinco estados $|2,m\rangle$. La proyección $L_z=m\hbar$ está definida, mientras que $L_x$ y $L_y$ quedan indeterminados (el vector yace en algún punto de su cono). Como $\sqrt{\ell(\ell+1)}>\ell$, ningún estado alinea $\vec L$ con el eje $z$.}
    \label{fig:cono_L}
    \label{fig:conos}
\end{figure}


Sin embargo, como demostramos en la Ec.~\eqref{eq:l2_comm}, el operador magnitud al cuadrado $\hat{L}^2$ conmuta con \textit{todas} las componentes individuales:
\begin{equation}
    [\hat{L}^2, \hat{L}_x] = 0, \quad [\hat{L}^2, \hat{L}_y] = 0, \quad [\hat{L}^2, \hat{L}_z] = 0.
\end{equation}
Esto nos permite elegir $\hat{L}^2$ y \textbf{una única componente espacial} (por convención universal, el eje polar $\hat{L}_z$) para formar un conjunto compatible:
\begin{equation}
    [\hat{L}^2, \hat{L}_z] = 0.
\end{equation}

\subsubsection{Base Ortonormal Común $\{\ket{\ell, m}\}$ y Operadores Compatibles}

El teorema fundamental de observables que conmutan garantiza que existe una base ortonormal simultánea común de autoestados para ambos operadores, denotada por $\{\ket{\ell, m}\}$:
\begin{align}
    \hat{L}^2 \ket{\ell, m} &= \hbar^2 \ell(\ell+1)\ket{\ell, m}, \quad \ell \in \{0, 1, 2, \dots\}, \label{eq:l2_eigenvalues} \\
    \hat{L}_z \ket{\ell, m} &= \hbar m \ket{\ell, m}, \quad m \in \{-\ell, -\ell+1, \dots, \ell-1, \ell\}. \label{eq:lz_eigenvalues}
\end{align}

\begin{teoremabox}[Definición: Conjunto Completo de Observables Compatibles (CCOC)]
    Un conjunto de operadores hermíticos $\{\hat{A}, \hat{B}, \hat{C}, \dots\}$ constituye un \textbf{CCOC} (o CSCO por sus siglas en inglés, \textit{Complete Set of Compatible Observables}) si:
    \begin{enumerate}
        \item Todos los operadores conmutan mutuamente entre sí por pares ($[\hat{A}, \hat{B}] = 0$, etc.).
        \item Existe una base ortonormal común única donde cada combinación de autovalores $\{a, b, c, \dots\}$ especifica \textbf{un único estado físico}, eliminando por completo cualquier degeneración espectral residual.
    \end{enumerate}
\end{teoremabox}

\subsubsection{Jerarquía del CCOC: Del Potencial Central al Espín del Electrón}

En un potencial esféricamente simétrico central $V(r)$ (como el átomo de hidrógeno), el hamiltoniano $\hat{H} = \frac{\hat{p}^2}{2m} + V(r)$ es invariante ante rotaciones espaciales, lo que implica:
\begin{equation}
    [\hat{H}, \hat{L}^2] = 0 \quad \text{y} \quad [\hat{H}, \hat{L}_z] = 0.
\end{equation}
Para los grados de libertad orbitales (espaciales), el \textbf{CCOC del potencial central} viene dado por:
\begin{equation}
    \boxed{\text{CCOC}_{\text{orbital}} = \{\hat{H}, \hat{L}^2, \hat{L}_z\}},
\end{equation}
cuyos autoestados estacionarios comunes están biunívocamente identificados por los números cuánticos orbitales $\{n, \ell, m_\ell\}$ mediante el estado $\ket{n, \ell, m_\ell}$ \cite{cohen, zettili}.\\
Cuando se incorpora el grado de libertad del \textbf{espín del electrón} ($s = 1/2$), como los operadores orbitales conmutan con los operadores de espín ($[\hat{\vec{L}}, \hat{\vec{S}}] = 0$), el CCOC completo del átomo se expande naturalmente a:
\begin{equation}
    \boxed{\text{CCOC}_{\text{completo}} = \{\hat{H}, \hat{L}^2, \hat{L}_z, \hat{S}^2, \hat{S}_z\}},
\end{equation}
especificando sin ambigüedad la base ortonormal completa $\ket{n, \ell, m_\ell, s, m_s}$, donde $m_s = \pm 1/2$. (Cabe notar que para un electrón individual con espín intrínseco fijo $s=1/2$, el operador Casimir $\hat{S}^2 = s(s+1)\hbar^2\hat{I} = \frac{3}{4}\hbar^2\hat{I}$ es un múltiplo escalar constante de la identidad, por lo que el observable dinámico que levanta la degeneración de espín es $\hat{S}_z$).

\subsection{Operadores Escalera de Momento Angular \texorpdfstring{$\hat{L}_\pm$}{L+-} y Cuantización Espectral}
\label{subsec:operadores_escalera}

Para determinar algebraicamente los autovalores y autoestados del momento angular sin resolver ecuaciones diferenciales en derivadas parciales, introducimos los \textbf{operadores escalera} (o de subida y bajada) \cite{sakurai, zettili}:
\begin{equation}
    \boxed{\hat{L}_+ \equiv \hat{L}_x + \iu\hat{L}_y, \qquad \hat{L}_- \equiv \hat{L}_x - \iu\hat{L}_y}.
    \label{eq:ladder_def}
\end{equation}
Como $\hat{L}_x$ y $\hat{L}_y$ son hermíticos, cada operador escalera es el adjunto hermítico del otro:
\begin{equation}
    (\hat{L}_+)\adj = (\hat{L}_x + \iu\hat{L}_y)\adj = \hat{L}_x\adj - \iu\hat{L}_y\adj = \hat{L}_x - \iu\hat{L}_y = \hat{L}_-.
\end{equation}
Inversamente, podemos recuperar las componentes cartesianas mediante las combinaciones:
\begin{equation}
    \hat{L}_x = \frac{1}{2}\left(\hat{L}_+ + \hat{L}_-\right), \qquad \hat{L}_y = \frac{1}{2\iu}\left(\hat{L}_+ - \hat{L}_-\right).
\end{equation}

\subsubsection{Álgebra Fundamental de los Operadores Escalera}

Utilizando el álgebra de Lie $[\hat{L}_i, \hat{L}_j] = \iu\hbar\varepsilon_{ijk}\hat{L}_k$, evaluamos los conmutadores de $\hat{L}_\pm$ con los operadores del CCOC:
\begin{enumerate}
    \item \textbf{Conmutador con la componente polar $\hat{L}_z$:}
    \begin{align}
        [\hat{L}_z, \hat{L}_\pm] &= [\hat{L}_z, \hat{L}_x \pm \iu\hat{L}_y] = [\hat{L}_z, \hat{L}_x] \pm \iu[\hat{L}_z, \hat{L}_y] \nonumber\\
        &= (\iu\hbar\hat{L}_y) \pm \iu(-\iu\hbar\hat{L}_x) = \iu\hbar\hat{L}_y \pm \hbar\hat{L}_x = \pm\hbar(\hat{L}_x \pm \iu\hat{L}_y) \nonumber\\
        &= \boxed{\pm\hbar\hat{L}_\pm}.
        \label{eq:comm_lz_lpm}
    \end{align}
    \item \textbf{Conmutador con el operador Casimir $\hat{L}^2$:}
    Como $[\hat{L}^2, \hat{L}_x] = 0$ y $[\hat{L}^2, \hat{L}_y] = 0$, se deduce inmediatamente por linealidad:
    \begin{equation}
        \boxed{[\hat{L}^2, \hat{L}_\pm] = 0}.
        \label{eq:comm_l2_lpm}
    \end{equation}
    \item \textbf{Conmutador mutuo entre subida y bajada:}
    \begin{align}
        [\hat{L}_+, \hat{L}_-] &= [\hat{L}_x + \iu\hat{L}_y, \, \hat{L}_x - \iu\hat{L}_y] = -\iu[\hat{L}_x, \hat{L}_y] + \iu[\hat{L}_y, \hat{L}_x] \nonumber\\
        &= -\iu(\iu\hbar\hat{L}_z) + \iu(-\iu\hbar\hat{L}_z) = \hbar\hat{L}_z + \hbar\hat{L}_z = \boxed{2\hbar\hat{L}_z}.
        \label{eq:comm_lplus_lminus}
    \end{align}
\end{enumerate}

\subsubsection{Descomposición de $\hat{L}^2$ y Acción sobre la Base $\{\ket{\ell, m}\}$}

Multiplicando los operadores escalera en ambos órdenes:
\begin{align}
    \hat{L}_-\hat{L}_+ &= (\hat{L}_x - \iu\hat{L}_y)(\hat{L}_x + \iu\hat{L}_y) = \hat{L}_x^2 + \hat{L}_y^2 + \iu[\hat{L}_x, \hat{L}_y] = \hat{L}_x^2 + \hat{L}_y^2 - \hbar\hat{L}_z, \\
    \hat{L}_+\hat{L}_- &= (\hat{L}_x + \iu\hat{L}_y)(\hat{L}_x - \iu\hat{L}_y) = \hat{L}_x^2 + \hat{L}_y^2 - \iu[\hat{L}_x, \hat{L}_y] = \hat{L}_x^2 + \hat{L}_y^2 + \hbar\hat{L}_z.
\end{align}
Sumando $\hat{L}_z^2$ a ambos lados y reconociendo $\hat{L}^2 = \hat{L}_x^2 + \hat{L}_y^2 + \hat{L}_z^2$, obtenemos las dos identidades maestras:
\begin{equation}
    \boxed{\hat{L}^2 = \hat{L}_\mp\hat{L}_\pm + \hat{L}_z^2 \pm \hbar\hat{L}_z = \frac{1}{2}(\hat{L}_+\hat{L}_- + \hat{L}_-\hat{L}_+) + \hat{L}_z^2}.
    \label{eq:l2_ladder_decomp}
\end{equation}
Examinemos ahora cómo actúa $\hat{L}_\pm$ sobre un autoestado común $\ket{\ell, m}$. Utilizando la Ec.~\eqref{eq:comm_lz_lpm}:
\begin{equation}
    \hat{L}_z (\hat{L}_\pm\ket{\ell, m}) = (\hat{L}_\pm\hat{L}_z \pm \hbar\hat{L}_\pm)\ket{\ell, m} = (\hbar m \pm \hbar)\hat{L}_\pm\ket{\ell, m} = \hbar(m \pm 1)(\hat{L}_\pm\ket{\ell, m}).
\end{equation}
¡El operador $\hat{L}_+$ incrementa el número cuántico magnético $m$ en una unidad (+1), mientras que $\hat{L}_-$ lo decrementa en una unidad (-1), preservando intacto el valor de $\ell$ debido a que $[\hat{L}^2, \hat{L}_\pm] = 0$!\\
Para calcular la constante de proporcionalidad (normalización), evaluamos la norma al cuadrado:
\begin{align}
    \norm{\hat{L}_\pm\ket{\ell, m}}^2 &= \bra{\ell, m}\hat{L}_\pm\adj\hat{L}_\pm\ket{\ell, m} = \bra{\ell, m}\hat{L}_\mp\hat{L}_\pm\ket{\ell, m} \nonumber\\
    &= \matrixel{\ell, m}{\hat{L}^2 - \hat{L}_z^2 \mp \hbar\hat{L}_z}{\ell, m} \nonumber\\
    &= \hbar^2\left[\ell(\ell+1) - m^2 \mp m\right] = \hbar^2\left[\ell(\ell+1) - m(m \pm 1)\right].
\end{align}
Eligiendo la convención de fase estándar de Condon--Shortley (coeficientes reales y positivos):
\begin{equation}
    \boxed{\hat{L}_\pm\ket{\ell, m} = \hbar\sqrt{\ell(\ell+1) - m(m \pm 1)}\,\ket{\ell, m \pm 1}}.
    \label{eq:ladder_action}
\end{equation}
En particular, para el valor máximo $m = +\ell$ y el mínimo $m = -\ell$:
\begin{equation}
    \hat{L}_+\ket{\ell, \ell} = 0, \qquad \hat{L}_-\ket{\ell, -\ell} = 0,
\end{equation}
lo que demuestra que el espectro de proyecciones está acotado en el intervalo finito $m \in \{-\ell, -\ell+1, \dots, \ell-1, \ell\}$, conteniendo exactamente $(2\ell+1)$ estados independientes.

\begin{figure}[H]
\centering
\begin{tikzpicture}[>=Stealth, font=\small, lev/.style={thick}]
  \def\h{1.15}   % separación vertical entre niveles

  % ---------------- l = 1 ----------------
  \begin{scope}[xshift=0cm]
    \node[font=\normalsize] at (0,3.75) {$\ell=1$};
    \foreach \m in {-1,0,1}{
      \draw[lev] (-1.05,\m*\h)--(1.05,\m*\h);
      \node[anchor=east] at (-1.1,\m*\h) {$m=\m$};
    }
    \foreach \m/\c in {-1/{\hbar\sqrt2},0/{\hbar\sqrt2}}{
      \draw[->,thick,azulUNI]   (-0.4,\m*\h+0.06)--(-0.4,\m*\h+\h-0.06);
      \draw[->,thick,naranjaUNI](0.4,\m*\h+\h-0.06)--(0.4,\m*\h+0.06);
      \node[azulUNI,anchor=east,font=\scriptsize] at (-0.48,\m*\h+0.5*\h) {$\c$};
    }
    \node[azulUNI,font=\scriptsize,align=center] at (0,1*\h+0.75)
         {$\hat L_+|1,1\rangle=0$};
    \node[naranjaUNI,font=\scriptsize,align=center] at (0,-1*\h-0.75)
         {$\hat L_-|1,-1\rangle=0$};
  \end{scope}

  % ---------------- l = 2 ----------------
  \begin{scope}[xshift=6.8cm]
    \node[font=\normalsize] at (0,3.75) {$\ell=2$};
    \foreach \m in {-2,-1,0,1,2}{
      \draw[lev] (-1.05,\m*\h)--(1.05,\m*\h);
      \node[anchor=east] at (-1.1,\m*\h) {$m=\m$};
    }
    \foreach \m/\c in {-2/{2\hbar},-1/{\hbar\sqrt6},0/{\hbar\sqrt6},1/{2\hbar}}{
      \draw[->,thick,azulUNI]   (-0.4,\m*\h+0.06)--(-0.4,\m*\h+\h-0.06);
      \draw[->,thick,naranjaUNI](0.4,\m*\h+\h-0.06)--(0.4,\m*\h+0.06);
      \node[azulUNI,anchor=east,font=\scriptsize] at (-0.48,\m*\h+0.5*\h) {$\c$};
    }
    \node[azulUNI,font=\scriptsize] at (0,2*\h+0.35)
         {$\hat L_+|2,2\rangle=0$};
    \node[naranjaUNI,font=\scriptsize] at (0,-2*\h-0.35)
         {$\hat L_-|2,-2\rangle=0$};
  \end{scope}

  % leyenda
  \draw[->,thick,azulUNI]    (2.7,-3.55)--(3.3,-3.55);
  \node[anchor=west,font=\scriptsize] at (3.35,-3.55) {$\hat L_+$ (sube)};
  \draw[->,thick,naranjaUNI] (5.1,-3.55)--(5.7,-3.55);
  \node[anchor=west,font=\scriptsize] at (5.75,-3.55) {$\hat L_-$ (baja)};
\end{tikzpicture}
\caption{Operadores escalera. Cada flecha sube o baja una unidad de $m$ sin
cambiar $\ell$. El coeficiente sobre la flecha es
$\hbar\sqrt{\ell(\ell+1)-m(m\pm1)}$ y es el mismo para $\hat L_+$ (de $m$ a
$m+1$) y para $\hat L_-$ (de $m+1$ a $m$). Para $\ell=1$ son las entradas
$\hbar\sqrt2$ de las matrices~\eqref{eq:matrix_ladder_1}.}
\label{fig:escalera}
\end{figure}


\subsection{Representación Matricial Explícita para \texorpdfstring{$\ell = 1$}{l = 1} (Espacio 3D)}
\label{subsec:matrices_l1}

Para el valor $\ell = 1$, la dimensión del subespacio es $2\ell + 1 = 3$. La base ortonormal estándar, ordenada según autovalores decrecientes de $\hat{L}_z$, es:
\begin{equation}
    \mathcal{B} = \left\{ \ket{1, 1} \equiv \begin{pmatrix} 1 \\ 0 \\ 0 \end{pmatrix}, \quad \ket{1, 0} \equiv \begin{pmatrix} 0 \\ 1 \\ 0 \end{pmatrix}, \quad \ket{1, -1} \equiv \begin{pmatrix} 0 \\ 0 \\ 1 \end{pmatrix} \right\}.
\end{equation}
Calculamos elemento a elemento las matrices representantes de cada operador en esta base:
\begin{enumerate}
    \item \textbf{Matriz de $\hat{L}_z$:} Puesto que la base es propia de $\hat{L}_z$ ($\hat{L}_z\ket{1, m} = m\hbar\ket{1, m}$), su matriz es puramente diagonal:
    \begin{equation}
        \hat{L}_z = \hbar \begin{pmatrix} 1 & 0 & 0 \\ 0 & 0 & 0 \\ 0 & 0 & -1 \end{pmatrix}.
        \label{eq:matrix_lz_1}
    \end{equation}
    \item \textbf{Matrices Escalera $\hat{L}_+$ y $\hat{L}_-$:} Aplicando la Ec.~\eqref{eq:ladder_action}:
    \begin{align}
        \hat{L}_+\ket{1, 1} &= 0, \\
        \hat{L}_+\ket{1, 0} &= \hbar\sqrt{1(2) - 0(1)}\ket{1, 1} = \hbar\sqrt{2}\ket{1, 1}, \\
        \hat{L}_+\ket{1, -1} &= \hbar\sqrt{1(2) - (-1)(0)}\ket{1, 0} = \hbar\sqrt{2}\ket{1, 0}.
    \end{align}
    Por ende, las matrices escalera de $3\times 3$ resultan ser:
    \begin{equation}
        \hat{L}_+ = \hbar\sqrt{2}\begin{pmatrix} 0 & 1 & 0 \\ 0 & 0 & 1 \\ 0 & 0 & 0 \end{pmatrix}, \qquad \hat{L}_- = (\hat{L}_+)\adj = \hbar\sqrt{2}\begin{pmatrix} 0 & 0 & 0 \\ 1 & 0 & 0 \\ 0 & 1 & 0 \end{pmatrix}.
        \label{eq:matrix_ladder_1}
    \end{equation}
    \item \textbf{Matrices Cartesianas $\hat{L}_x$ y $\hat{L}_y$:}
    \begin{align}
        \hat{L}_x &= \frac{1}{2}(\hat{L}_+ + \hat{L}_-) = \frac{\hbar}{\sqrt{2}}\begin{pmatrix} 0 & 1 & 0 \\ 1 & 0 & 1 \\ 0 & 1 & 0 \end{pmatrix}, \label{eq:matrix_lx_1} \\
        \hat{L}_y &= \frac{1}{2\iu}(\hat{L}_+ - \hat{L}_-) = \frac{\hbar}{\sqrt{2}}\begin{pmatrix} 0 & -\iu & 0 \\ \iu & 0 & -\iu \\ 0 & \iu & 0 \end{pmatrix}. \label{eq:matrix_ly_1}
    \end{align}
    \item \textbf{Matriz del Casimir $\hat{L}^2$:}
    Puesto que $\hat{L}^2\ket{1, m} = \hbar^2 1(1+1)\ket{1, m} = 2\hbar^2\ket{1, m}$, es proporcional a la identidad $3\times 3$:
    \begin{equation}
        \hat{L}^2 = 2\hbar^2 \begin{pmatrix} 1 & 0 & 0 \\ 0 & 1 & 0 \\ 0 & 0 & 1 \end{pmatrix} = 2\hbar^2 \hat{I}_3.
    \end{equation}
\end{enumerate}

\begin{alertabox}[Verificación Matricial Directa del Álgebra de Conmutadores]
    Comprobemos matricialmente el conmutador $[\hat{L}_x, \hat{L}_y]$:
    \begin{align}
        \hat{L}_x\hat{L}_y &= \frac{\hbar^2}{2}\begin{pmatrix} 0 & 1 & 0 \\ 1 & 0 & 1 \\ 0 & 1 & 0 \end{pmatrix}\begin{pmatrix} 0 & -\iu & 0 \\ \iu & 0 & -\iu \\ 0 & \iu & 0 \end{pmatrix} = \frac{\hbar^2}{2}\begin{pmatrix} \iu & 0 & -\iu \\ 0 & 0 & 0 \\ \iu & 0 & -\iu \end{pmatrix}, \\
        \hat{L}_y\hat{L}_x &= \frac{\hbar^2}{2}\begin{pmatrix} 0 & -\iu & 0 \\ \iu & 0 & -\iu \\ 0 & \iu & 0 \end{pmatrix}\begin{pmatrix} 0 & 1 & 0 \\ 1 & 0 & 1 \\ 0 & 1 & 0 \end{pmatrix} = \frac{\hbar^2}{2}\begin{pmatrix} -\iu & 0 & -\iu \\ 0 & 0 & 0 \\ \iu & 0 & \iu \end{pmatrix}.
    \end{align}
    Restando ambas matrices:
    \begin{equation}
        [\hat{L}_x, \hat{L}_y] = \frac{\hbar^2}{2}\begin{pmatrix} 2\iu & 0 & 0 \\ 0 & 0 & 0 \\ 0 & 0 & -2\iu \end{pmatrix} = \iu\hbar\left[ \hbar\begin{pmatrix} 1 & 0 & 0 \\ 0 & 0 & 0 \\ 0 & 0 & -1 \end{pmatrix} \right] = \iu\hbar\hat{L}_z.
    \end{equation}
    ¡El álgebra de Lie se satisface con total precisión aritmética en la representación matricial tridimensional!
\end{alertabox}

\subsection{El Hamiltoniano Central, el Laplaciano Esférico y la Separación de Variables}
\label{subsec:hamiltoniano_central_laplaciano}

En coordenadas esféricas $(r, \theta, \phi)$, el operador laplaciano $\nabla^2$ se descompone naturalmente en una parte radial y una parte angular:
\begin{equation}
    \nabla^2 = \frac{1}{r^2}\frac{\partial}{\partial r}\left(r^2\frac{\partial}{\partial r}\right) + \frac{1}{r^2}\left[ \frac{1}{\sin\theta}\frac{\partial}{\partial\theta}\left(\sin\theta\frac{\partial}{\partial\theta}\right) + \frac{1}{\sin^2\theta}\frac{\partial^2}{\partial\phi^2} \right].
    \label{eq:laplaciano_esferico}
\end{equation}
Reconociendo la representación diferencial del operador Casimir $\hat{L}^2$ en coordenadas esféricas:
\begin{equation}
    \hat{L}^2 = -\hbar^2\left[ \frac{1}{\sin\theta}\frac{\partial}{\partial\theta}\left(\sin\theta\frac{\partial}{\partial\theta}\right) + \frac{1}{\sin^2\theta}\frac{\partial^2}{\partial\phi^2} \right], \qquad \hat{L}_z = -\iu\hbar\frac{\partial}{\partial\phi}.
\end{equation}
el Laplaciano se reescribe de forma sumamente compacta y elegante como:
\begin{equation}
    \boxed{\nabla^2 = \frac{1}{r^2}\frac{\partial}{\partial r}\left(r^2\frac{\partial}{\partial r}\right) - \frac{\hat{L}^2}{\hbar^2 r^2}}.
    \label{eq:laplaciano_con_L2}
\end{equation}
Sustituyendo esto en el Hamiltoniano de una partícula de masa reducida $\mu$ sometida a un potencial esférico central $V(r)$:
\begin{equation}
    \boxed{\hat{H} = -\frac{\hbar^2}{2\mu}\frac{1}{r^2}\frac{\partial}{\partial r}\left(r^2\frac{\partial}{\partial r}\right) + \frac{\hat{L}^2}{2\mu r^2} + V(r)}.
    \label{eq:hamiltoniano_central_expandido}
\end{equation}

\begin{teoremabox}[Invariancia y Conmutación en Potenciales Centrales]
    Para cualquier potencial con simetría esférica $V(r)$, el Hamiltoniano conmuta con $\hat{L}^2$ y con todas las componentes $\hat{L}_i$:
    \begin{equation}
        [\hat{H}, \hat{L}^2] = 0, \qquad [\hat{H}, \hat{L}_z] = 0.
    \end{equation}
\end{teoremabox}

\textbf{Demostración Didáctica:}\\
1. Los operadores radiales $\frac{\partial}{\partial r}$ y las funciones del radio $r$ y $V(r)$ actúan exclusivamente sobre la coordenada $r$, por lo que conmutan idénticamente con los operadores diferenciales angulares $\hat{L}^2$ y $\hat{L}_z$ que sólo actúan sobre $(\theta, \phi)$.\\
2. Por otro lado, la única dependencia angular del Hamiltoniano $\hat{H}$ en la Ec.~\eqref{eq:hamiltoniano_central_expandido} está contenida en el término $\frac{\hat{L}^2}{2\mu r^2}$. Como $[\hat{L}^2, \hat{L}^2] = 0$ y $[\hat{L}^2, \hat{L}_z] = 0$ (demostrado en la Ec.~\eqref{eq:l2_comm}), concluimos de inmediato que:
\begin{equation}
    [\hat{H}, \hat{L}^2] = 0 \quad \text{y} \quad [\hat{H}, \hat{L}_z] = 0.
\end{equation}
El Teorema 2.3 garantiza la existencia de una base ortonormal común de autoestados simultáneos para el CCOC $\{\hat{H}, \hat{L}^2, \hat{L}_z\}$. Debido a la estructura diferencial del laplaciano en coordenadas esféricas (donde la parte angular coincide con $-\hat{L}^2/(\hbar^2 r^2)$), el método de separación de variables conduce de forma natural a la factorización espacial de las funciones de onda:
\begin{equation}
    \psi_{n\ell m}(r, \theta, \phi) = R_{n\ell}(r)\,Y_\ell^m(\theta, \phi),
\end{equation}
donde los \textbf{armónicos esféricos} $Y_\ell^m(\theta, \phi)$ son los autoestados comunes de $\hat{L}^2$ y $\hat{L}_z$, y la función radial $R_{n\ell}(r)$ satisface la ecuación diferencial 1D:
\begin{equation}
    \left[ -\frac{\hbar^2}{2\mu}\frac{1}{r^2}\frac{\dif}{\dif r}\left(r^2\frac{\dif}{\dif r}\right) + \frac{\hbar^2\ell(\ell+1)}{2\mu r^2} + V(r) \right] R_{n\ell}(r) = E_{n\ell} R_{n\ell}(r).
\end{equation}
El término $\frac{\hbar^2\ell(\ell+1)}{2\mu r^2}$ representa la \textbf{barrera centrífuga cuántica}: una fuerza repulsiva de origen cinemático que aleja a la partícula del origen para todo estado con momento angular no nulo ($\ell \ge 1$).

\subsection{El Producto Tensorial \texorpdfstring{$\mathcal{H}_1 \otimes \mathcal{H}_2$}{H1 x H2} y el Acoplamiento de Momentos Angulares}
\label{subsec:producto_tensorial_acoplamiento}

Cuando consideramos sistemas físicos compuestos por múltiples partículas (por ejemplo, dos electrones) o grados de libertad independientes de una misma partícula (como el movimiento espacial orbital acoplado al espín intrínseco), el espacio de estados global se construye mediante el \textbf{producto tensorial de espacios de Hilbert} \cite{sakurai, cohen}.

\begin{teoremabox}[Definición Formal: Espacio Producto Tensorial $\mathcal{H}_1 \otimes \mathcal{H}_2$]
    Sean $\mathcal{H}_1$ y $\mathcal{H}_2$ dos espacios de Hilbert de dimensiones $N_1$ y $N_2$ con bases ortonormales $\{\ket{u_i}\}_{i=1}^{N_1}$ y $\{\ket{v_j}\}_{j=1}^{N_2}$ respectivamente. El espacio producto tensorial $\mathcal{H} = \mathcal{H}_1 \otimes \mathcal{H}_2$ satisface:
    \begin{enumerate}
        \item \textbf{Bilinealidad:} Para todo $\alpha, \beta \in \mathbb{C}$:
        \begin{equation}
            (\alpha\ket{u_1} + \beta\ket{u_2}) \otimes \ket{v} = \alpha(\ket{u_1} \otimes \ket{v}) + \beta(\ket{u_2} \otimes \ket{v}).
        \end{equation}
        \item \textbf{Producto Interno Compuesto:} Para estados producto $\ket{\Phi} = \ket{\phi_1} \otimes \ket{\phi_2}$ y $\ket{\Psi} = \ket{\psi_1} \otimes \ket{\psi_2}$:
        \begin{equation}
            \braket{\Phi}{\Psi} = \braket{\phi_1}{\psi_1}_{\mathcal{H}_1} \cdot \braket{\phi_2}{\psi_2}_{\mathcal{H}_2}.
        \end{equation}
        \item \textbf{Dimensión Total:} La dimensión es el producto de las dimensiones:
        \begin{equation}
            \dim(\mathcal{H}_1 \otimes \mathcal{H}_2) = \dim(\mathcal{H}_1) \cdot \dim(\mathcal{H}_2) = N_1 \cdot N_2.
        \end{equation}
        \item \textbf{Base Producto Desacoplada:} El conjunto de pares ordenados $\{\ket{u_i} \otimes \ket{v_j}\} \equiv \{\ket{u_i, v_j}\}$ conforma una base ortonormal completa de $\mathcal{H}_1 \otimes \mathcal{H}_2$, satisfaciendo la relación de clausura:
        \begin{equation}
            \sum_{i=1}^{N_1} \sum_{j=1}^{N_2} \ket{u_i, v_j}\bra{u_i, v_j} = \hat{I}_1 \otimes \hat{I}_2 = \hat{I}.
        \end{equation}
        \item \textbf{Producto de Operadores Compuestos (Propiedad Fundamental):}
        \begin{equation}
            \boxed{(\hat{A}_1 \otimes \hat{B}_1)(\hat{A}_2 \otimes \hat{B}_2) = (\hat{A}_1\hat{A}_2) \otimes (\hat{B}_1\hat{B}_2)}.
            \label{eq:operator_tensor_product_rule}
        \end{equation}
    \end{enumerate}
\end{teoremabox}

\subsubsection{Motivación y Necesidad Física: El Acoplamiento Espín-Órbita ($LS$)}
\label{subsec:acoplamiento_espin_orbita}

¿Por qué resulta indispensable en la física cuántica real sumar momentos angulares y abandonar la descripción individual de cada subsistema?

En un átomo idealizado sin interacciones magnéticas relativistas internas, el Hamiltoniano no perturbado $\hat{H}_0$ conmuta independientemente tanto con el momento angular orbital $\hat{\vec{L}}$ como con el espín intrínseco $\hat{\vec{S}}$:
\begin{equation}
    [\hat{H}_0, \hat{\vec{L}}] = 0 \quad \text{y} \quad [\hat{H}_0, \hat{\vec{S}}] = 0.
\end{equation}
En dicho régimen simplificado, las componentes individuales $\hat{L}_z$ y $\hat{S}_z$ se conservan de forma separada en el tiempo.

Sin embargo, en el mundo real, un electrón que orbita a velocidades relativistas alrededor del núcleo experimenta en su propio sistema de referencia en reposo un intenso campo magnético efectivo generado por el movimiento aparente del núcleo cargado. Dicho campo magnético interno interactúa con el momento magnético dipolar intrínseco del electrón, generando el término de \textbf{acoplamiento espín-órbita}:
\begin{equation}
    \hat{H}_{SO} = \xi(r)\,\hat{\vec{L}}\cdot\hat{\vec{S}} = \xi(r)\left( \hat{L}_x\hat{S}_x + \hat{L}_y\hat{S}_y + \hat{L}_z\hat{S}_z \right), \quad \xi(r) > 0.
    \label{eq:h_so_def}
\end{equation}
Al evaluar la conmutación de este término con las componentes orbitales individuales:
\begin{align}
    [\hat{H}_{SO}, \hat{L}_i] &= \xi(r)[\hat{L}_j\hat{S}_j, \, \hat{L}_i] = \xi(r)[\hat{L}_j, \hat{L}_i]\hat{S}_j = \xi(r)(-\iu\hbar\,\varepsilon_{jik}\hat{L}_k)\hat{S}_j \nonumber\\
    &= \iu\hbar\,\xi(r)\,\varepsilon_{ijk}\hat{L}_j\hat{S}_k \neq 0.
\end{align}
Análogamente, para el espín: $[\hat{H}_{SO}, \hat{S}_i] = -\iu\hbar\,\xi(r)\,\varepsilon_{ijk}\hat{L}_j\hat{S}_k \neq 0$.

\begin{alertabox}[Consecuencia Física Fundamental: Ruptura de Conservación Aislada]
    Debido a la interacción espín-órbita ($\propto \hat{\vec{L}}\cdot\hat{\vec{S}}$), el momento angular orbital $\hat{\vec{L}}$ y el espín $\hat{\vec{S}}$ \textbf{dejan de conmutar individualmente con el Hamiltoniano total}:
    \begin{equation}
        [\hat{H}_{\text{total}}, \hat{\vec{L}}] \neq 0 \quad \text{y} \quad [\hat{H}_{\text{total}}, \hat{\vec{S}}] \neq 0.
    \end{equation}
    Por consiguiente, por el Teorema de Ehrenfest:
    \begin{equation}
        \frac{\dif}{\dif t}\ev{\hat{\vec{L}}} \neq 0, \qquad \frac{\dif}{\dif t}\ev{\hat{\vec{S}}} \neq 0.
    \end{equation}
    \textbf{¡Ni el momento orbital ni el espín se conservan de manera aislada!} Un torque magnético interno transfiere momento angular de forma continua entre el movimiento orbital y el espín.
\end{alertabox}

\subsubsection{Suma de Momentos Angulares y el Operador Momento Angular Total $\hat{\vec{J}}$}

Sin embargo, al sumar ambas variaciones temporales:
\begin{equation}
    [\hat{H}_{SO}, \, \hat{L}_i + \hat{S}_i] = \iu\hbar\,\xi(r)\,\varepsilon_{ijk}\hat{L}_j\hat{S}_k - \iu\hbar\,\xi(r)\,\varepsilon_{ijk}\hat{L}_j\hat{S}_k = 0.
\end{equation}
¡La suma total conmuta exactamente con el Hamiltoniano! Esto motiva la definición del \textbf{operador momento angular total}:
\begin{equation}
    \boxed{\hat{\vec{J}} \equiv \hat{\vec{L}} + \hat{\vec{S}} = \hat{\vec{L}} \otimes \hat{I}_S + \hat{I}_L \otimes \hat{\vec{S}}}, \qquad [\hat{H}_{\text{total}}, \hat{\vec{J}}] = 0.
\end{equation}
En un contexto más general donde se acoplan dos subsistemas cualesquiera (dos electrones, o dos grados de libertad independientes) con operadores $\hat{\vec{J}}_1$ y $\hat{\vec{J}}_2$:
\begin{equation}
    [\hat{J}_{1i}, \hat{J}_{1j}] = \iu\hbar\,\varepsilon_{ijk}\hat{J}_{1k}, \qquad [\hat{J}_{2i}, \hat{J}_{2j}] = \iu\hbar\,\varepsilon_{ijk}\hat{J}_{2k}, \qquad [\hat{J}_{1i} \otimes \hat{I}_2, \, \hat{I}_1 \otimes \hat{J}_{2j}] = 0.
\end{equation}
Definimos el \textbf{operador momento angular total} como:
\begin{equation}
    \boxed{\hat{\vec{J}} \equiv \hat{\vec{J}}_1 \otimes \hat{I}_2 + \hat{I}_1 \otimes \hat{\vec{J}}_2 \equiv \hat{\vec{J}}_1 + \hat{\vec{J}}_2}.
\end{equation}
Comprobemos que el operador suma satisface el álgebra de Lie del momento angular empleando la regla del producto tensorial \eqref{eq:operator_tensor_product_rule}:
\begin{align}
    [\hat{J}_i, \hat{J}_j] &= [\hat{J}_{1i} + \hat{J}_{2i}, \, \hat{J}_{1j} + \hat{J}_{2j}] = [\hat{J}_{1i}, \hat{J}_{1j}] + [\hat{J}_{2i}, \hat{J}_{2j}] + \underbrace{[\hat{J}_{1i}, \hat{J}_{2j}]}_{0} + \underbrace{[\hat{J}_{2i}, \hat{J}_{1j}]}_{0} \nonumber\\
    &= \iu\hbar\,\varepsilon_{ijk}\hat{J}_{1k} + \iu\hbar\,\varepsilon_{ijk}\hat{J}_{2k} = \iu\hbar\,\varepsilon_{ijk}(\hat{J}_{1k} + \hat{J}_{2k}) = \boxed{\iu\hbar\,\varepsilon_{ijk}\hat{J}_k}.
\end{align}
¡El operador total $\hat{\vec{J}}$ es un momento angular genuino y riguroso!

\subsubsection{Incompatibilidad entre la Base Desacoplada y la Base Acoplada}

Consideremos el operador Casimir total al cuadrado:
\begin{equation}
    \hat{J}^2 = (\hat{\vec{J}}_1 + \hat{\vec{J}}_2)^2 = \hat{J}_1^2 + \hat{J}_2^2 + 2\hat{\vec{J}}_1\cdot\hat{\vec{J}}_2 = \hat{J}_1^2 + \hat{J}_2^2 + 2\hat{J}_{1z}\hat{J}_{2z} + \hat{J}_{1+}\hat{J}_{2-} + \hat{J}_{1-}\hat{J}_{2+}.
    \label{eq:j2_expansion_coupled}
\end{equation}
Evaluemos las relaciones de compatibilidad de este operador con los operadores individuales:
\begin{itemize}
    \item Conmuta con las magnitudes individuales al cuadrado: $[\hat{J}^2, \hat{J}_1^2] = 0$ y $[\hat{J}^2, \hat{J}_2^2] = 0$.
    \item Conmuta con la componente polar total: $[\hat{J}^2, \hat{J}_z] = 0$, donde $\hat{J}_z = \hat{J}_{1z} + \hat{J}_{2z}$.
    \item \textbf{¡NO conmuta con las componentes polares individuales $\hat{J}_{1z}$ y $\hat{J}_{2z}$!}\\
    En efecto, calculando el conmutador con $\hat{J}_{1z}$:
    \begin{align}
        [\hat{J}^2, \hat{J}_{1z}] &= [\hat{J}_1^2 + \hat{J}_2^2 + 2(\hat{J}_{1x}\hat{J}_{2x} + \hat{J}_{1y}\hat{J}_{2y} + \hat{J}_{1z}\hat{J}_{2z}), \, \hat{J}_{1z}] \nonumber\\
        &= 2[\hat{J}_{1x}, \hat{J}_{1z}]\hat{J}_{2x} + 2[\hat{J}_{1y}, \hat{J}_{1z}]\hat{J}_{2y} = 2(-\iu\hbar\hat{J}_{1y})\hat{J}_{2x} + 2(\iu\hbar\hat{J}_{1x})\hat{J}_{2y} \nonumber\\
        &= \boxed{2\iu\hbar\left(\hat{J}_{1x}\hat{J}_{2y} - \hat{J}_{1y}\hat{J}_{2x}\right) \neq 0 \quad (\text{en general})}.
    \end{align}
\end{itemize}
Esta no conmutatividad es el origen de los dos esquemas de descripción cuántica de un sistema compuesto:
\begin{enumerate}
    \item \textbf{Base Desacoplada $\{\ket{j_1, m_1; j_2, m_2}\}$:} Es la base simultánea del CCOC desacoplado $\{\hat{J}_1^2, \hat{J}_2^2, \hat{J}_{1z}, \hat{J}_{2z}\}$. En esta base, las proyecciones individuales $m_1$ y $m_2$ están perfectamente definidas, pero la magnitud total del momento angular $J$ está completamente indeterminada.
    \item \textbf{Base Acoplada $\{\ket{j, M; j_1, j_2}\}$:} Es la base simultánea del CCOC acoplado $\{\hat{J}_1^2, \hat{J}_2^2, \hat{J}^2, \hat{J}_z\}$. En esta base, la magnitud total $J$ y la proyección total $M = m_1 + m_2$ son nítidas, pero las componentes individuales $m_1$ y $m_2$ no tienen valor definido.
\end{enumerate}

\subsubsection{Regla de Suma de Momentos Angulares y Coeficientes de Clebsch--Gordan}

Para dos momentos angulares dados $j_1$ y $j_2$, los valores permitidos para el momento angular total $j$ recorren el intervalo dado por la \textbf{desigualdad triangular cuántica}:
\begin{equation}
    \boxed{|j_1 - j_2| \le j \le j_1 + j_2}, \quad \text{con paso entero } \Delta j = 1.
\end{equation}
El espacio producto total se descompone en suma directa de subespacios irreductibles:
\begin{equation}
    \mathcal{H}_{j_1} \otimes \mathcal{H}_{j_2} = \bigoplus_{j=|j_1-j_2|}^{j_1+j_2} \mathcal{H}_j.
\end{equation}
Comprobación de la conservación de la dimensión total:
\begin{equation}
    \sum_{j=|j_1-j_2|}^{j_1+j_2} (2j + 1) = (2j_1 + 1)(2j_2 + 1).
\end{equation}
El cambio de base unitario entre la base desacoplada y la acoplada se realiza mediante los \textbf{Coeficientes de Clebsch--Gordan} $C_{m_1, m_2; M}^{j_1, j_2; j} \equiv \braket{j_1, m_1; j_2, m_2}{j, M}$:
\begin{equation}
    \boxed{\ket{j, M} = \sum_{m_1=-j_1}^{j_1} \sum_{m_2=-j_2}^{j_2} \braket{j_1, m_1; j_2, m_2}{j, M} \ket{j_1, m_1; j_2, m_2}}.
\end{equation}
Un coeficiente de Clebsch--Gordan es estrictamente nulo a menos que satisfaga la regla de conservación de la proyección magnética:
\begin{equation}
    \braket{j_1, m_1; j_2, m_2}{j, M} = 0 \quad \text{si } M \neq m_1 + m_2.
\end{equation}

\subsubsection{Ejemplo Canónico 1: Acoplamiento de Dos Espines \texorpdfstring{$s_1 = 1/2$}{s1=1/2} y \texorpdfstring{$s_2 = 1/2$}{s2=1/2}}

La base desacoplada $\{\ket{++}, \ket{+-}, \ket{-+}, \ket{--}\}$ tiene dimensión $2\times 2 = 4$. Los valores posibles de $S$ son:
\begin{equation}
    S \in \{|1/2 - 1/2|, \, 1/2 + 1/2\} = \{0, 1\}.
\end{equation}
\begin{itemize}
    \item \textbf{Estado Triplete ($S = 1$, multiplicidad $2S+1 = 3$):}
    \begin{align}
        \ket{S=1, M=1} &= \ket{++}, \\
        \ket{S=1, M=0} &= \frac{1}{\sqrt{2}}\left( \ket{+-} + \ket{-+} \right), \\
        \ket{S=1, M=-1} &= \ket{--},
    \end{align}
    donde el estado con $M = 0$ se obtiene aplicando el operador de descenso $\hat{S}_- = \hat{S}_{1-} + \hat{S}_{2-}$ sobre el estado de máxima proyección $\ket{1, 1}$.
    \item \textbf{Estado Singlete ($S = 0$, multiplicidad $2S+1 = 1$):}
    Por ortogonalidad con el estado triplete $\ket{1, 0}$ en el mismo subespacio $M = 0$:
    \begin{equation}
        \ket{S=0, M=0} = \frac{1}{\sqrt{2}}\left( \ket{+-} - \ket{-+} \right).
    \end{equation}
\end{itemize}
¡El estado singlete es estrictamente antisimétrico ante el intercambio de las dos partículas y representa el célebre estado de entrelazamiento cuántico máximo (estado de Bell)!

\subsubsection{Ejemplo Canónico 2: Acoplamiento \texorpdfstring{$\ell_1 = 1 \otimes \ell_2 = 1$}{l1=1 x l2=1} y Matriz del Bloque \texorpdfstring{$M = 0$}{M=0}}

Consideremos dos momentos angulares de valor $\ell_1 = 1$ y $\ell_2 = 1$. La dimensión total del espacio es $(2\cdot 1 + 1)(2\cdot 1 + 1) = 3 \times 3 = 9$. Los multipletes totales posibles son:
\begin{equation}
    j \in \{0, 1, 2\}, \qquad \dim = 5 + 3 + 1 = 9.
\end{equation}
Para verificar algebraicamente que el operador Casimir $\hat{J}^2$ no es diagonal en la base desacoplada $\{\ket{m_1, m_2}\}$, utilicemos su expresión operativa \eqref{eq:j2_expansion_coupled}:
\begin{equation}
    \hat{J}^2 = \hat{L}_1^2 + \hat{L}_2^2 + 2\hat{L}_{1z}\hat{L}_{2z} + \hat{L}_{1+}\hat{L}_{2-} + \hat{L}_{1-}\hat{L}_{2+}.
\end{equation}
Evaluemos el bloque con proyección magnética nula $M = m_1 + m_2 = 0$. En la base desacoplada ordenada $\{\ket{1, -1}, \ket{0, 0}, \ket{-1, 1}\}$:
\begin{align}
    \hat{J}^2\ket{1, -1} &= \left(2\hbar^2 + 2\hbar^2 + 2\hbar^2(1)(-1)\right)\ket{1, -1} + \hat{L}_{1-}\hat{L}_{2+}\ket{1, -1} \nonumber\\
    &= 2\hbar^2\ket{1, -1} + (\hat{L}_{1-}\ket{1}_1) \otimes (\hat{L}_{2+}\ket{-1}_2) \nonumber\\
    &= 2\hbar^2\ket{1, -1} + (\hbar\sqrt{2}\ket{0}_1) \otimes (\hbar\sqrt{2}\ket{0}_2) = 2\hbar^2\ket{1, -1} + 2\hbar^2\ket{0, 0}.
\end{align}
¡Aparece un elemento fuera de la diagonal proporcional a $2\hbar^2$ que acopla $\ket{1, -1}$ con $\ket{0, 0}$!\\
Efectuando el mismo cálculo sobre los otros dos vectores, la matriz representante de $\hat{J}^2$ en el subespacio $M = 0$ resulta ser exactamente:
\begin{equation}
    \hat{J}^2|_{M=0} = \hbar^2 \begin{pmatrix} 2 & 2 & 0 \\ 2 & 4 & 2 \\ 0 & 2 & 2 \end{pmatrix}.
    \label{eq:matrix_j2_block_zero}
\end{equation}
Calculando los autovalores de esta matriz de $3\times 3$:
\begin{equation}
    \det\left( \hat{J}^2|_{M=0} - \lambda I \right) = -\lambda(\lambda - 2\hbar^2)(\lambda - 6\hbar^2) = 0 \implies \lambda \in \{0, \, 2\hbar^2, \, 6\hbar^2\}.
\end{equation}
Reconocemos de inmediato los valores de $\hbar^2 j(j+1)$:
\begin{itemize}
    \item Para $j = 0$: $\hbar^2 0(1) = 0$.
    \item Para $j = 1$: $\hbar^2 1(2) = 2\hbar^2$.
    \item Para $j = 2$: $\hbar^2 2(3) = 6\hbar^2$.
\end{itemize}
La traza de la matriz es $2 + 4 + 2 = 8 = 0 + 2 + 6$, conservada de forma idéntica. ¡En la base acoplada $\{\ket{j, M}\}$, la matriz $\hat{J}^2$ es estrictamente diagonal!

\begin{figure}[H]
\centering
\begin{tikzpicture}[x=2.2cm,y=1.6cm,>=Stealth,font=\small,
    ket/.style={rounded corners=3pt,draw=grisUNI!70,fill=white,
                minimum width=1.25cm,minimum height=0.7cm,inner sep=2pt},
    cero/.style={ket,draw=azulUNI,fill=azulUNI!15,very thick}]

  % rótulos de ejes
  \node[font=\footnotesize] at (0,-1.95) {$m_1\ (\text{partícula }1)$};
  \node[font=\footnotesize,rotate=90] at (-1.75,0) {$m_2\ (\text{partícula }2)$};

  % diagonales M = m1 + m2 = const
  \draw[dashed,grisUNI!70] (0.65,1.35)--(1.35,0.65)   node[right,font=\footnotesize]{$M=2$};
  \draw[dashed,grisUNI!70] (-0.35,1.35)--(1.35,-0.35) node[right,font=\footnotesize]{$M=1$};
  \draw[dashed,azulUNI]    (-1.35,1.35)--(1.35,-1.35) node[right,font=\footnotesize,azulUNI]{$M=0$};
  \draw[dashed,grisUNI!70] (-1.35,0.35)--(0.35,-1.35) node[right,font=\footnotesize]{$M=-1$};
  \draw[dashed,grisUNI!70] (-1.35,-0.65)--(-0.65,-1.35) node[right,font=\footnotesize]{$M=-2$};

  % estados |m1,m2>
  \foreach \a in {-1,0,1}{
    \foreach \b in {-1,0,1}{
      \pgfmathtruncatemacro{\s}{\a+\b}
      \ifnum\s=0
        \node[cero] at (\a,\b) {$|{\a},{\b}\rangle$};
      \else
        \node[ket]  at (\a,\b) {$|{\a},{\b}\rangle$};
      \fi
    }
  }

  % acoplamientos de J^2 dentro de M = 0
  \draw[<->,very thick,rojoUNI] (0.77,-0.77)--(0.23,-0.23)
        node[midway,above right,font=\footnotesize]{$2\hbar^2$};
  \draw[<->,very thick,rojoUNI] (-0.23,0.23)--(-0.77,0.77)
        node[midway,above right,font=\footnotesize]{$2\hbar^2$};

  % matriz del bloque M = 0
  \node[anchor=west,align=left] at (canvas cs:x=4.2cm,y=0.0cm) {%
    $\hat J^{2}\big|_{M=0}=\hbar^2
     \begin{pmatrix}2&2&0\\2&4&2\\0&2&2\end{pmatrix}$\\[8pt]
    \footnotesize autovalores $\{0,\,2\hbar^2,\,6\hbar^2\}$\\
    \footnotesize $\Rightarrow j=0,1,2$\\[6pt]
    \footnotesize $\dim=5+3+1=9$};
\end{tikzpicture}
\caption{Acoplamiento $\ell_1=1\otimes\ell_2=1$ en la base desacoplada
$|m_1,m_2\rangle$. $\hat J^2$ conserva $M=m_1+m_2$, por lo que sólo mezcla
estados sobre la misma diagonal discontinua. En el bloque $M=0$ (azul) los
términos $\hat L_{1\pm}\hat L_{2\mp}$ generan los elementos fuera de la
diagonal $2\hbar^2$ de la ec.~\eqref{eq:matrix_j2_block_zero}.}
\label{fig:cuadricula-m1m2}
\end{figure}

